\hypertarget{appendix-db3aef8f2fd6}{}\label{referencestart:appendix-db3aef8f2fd6}
\section{Cubic Hermite interpolation}\label{appendix-db3aef8f2fd6-cubic-hermite-interpolation}

\hypertarget{appendix-db3aef8f2fd6-cubic-hermite-interpolation}{}

\textbf{Curriculum depth:} Developed · \textbf{First seen:} Phase 10 · \textbf{Latest development:} Phase 10

\begin{allmightroledefinition}

\subsection{Definition and mechanism}\label{appendix-db3aef8f2fd6-definition-and-mechanism}

\hypertarget{appendix-db3aef8f2fd6-definition-and-mechanism}{}

A cubic Hermite segment is determined by endpoint values and endpoint tangent data. With normalized \(t\in[0,1]\), its four basis polynomials blend the two endpoint positions and two tangents. Tangents are derivatives with respect to the parameter, so physical-time derivatives must be scaled by segment duration: a velocity in units per second is not interchangeable with a tangent per normalized parameter. Hermite curves are useful when a camera path or animation must pass through known positions while exposing direct control over its entry and exit directions. The endpoint constraints determine one cubic polynomial; continuity between adjacent segments still requires compatible endpoint tangents.

\end{allmightroledefinition}

\subsection{Prerequisites}\label{appendix-db3aef8f2fd6-prerequisites}

\hypertarget{appendix-db3aef8f2fd6-prerequisites}{}

\begin{itemize}
\tightlist
\item
  \hyperref[appendix-0c2a5f69851e]{Interpolation versus approximation}
\item
  \href{../03-Geometry-Cameras-and-Image-Formation-P01/main.pdf\#appendix-f28eeec97a3d}{Spline bases and curve evaluation (Geometry, Cameras, and Image Formation, Part 1, Reference notes: Spline bases and curve evaluation)}
\end{itemize}

\subsection{Used by}\label{appendix-db3aef8f2fd6-used-by}

\hypertarget{appendix-db3aef8f2fd6-used-by}{}

\begin{itemize}
\tightlist
\item
  No downstream concept node is added in this stage.
\end{itemize}

\subsection{Full teaching explanation}\label{appendix-db3aef8f2fd6-full-teaching-explanation}

\hypertarget{appendix-db3aef8f2fd6-full-teaching-explanation}{}

\begin{itemize}
\tightlist
\item
  \href{../02-Mathematics-for-Visual-Computing-P02/main.pdf\#02-foundations-ii-ch028-reconstructing-values-and-curves-between-samples}{Reconstructing values and curves between samples (Mathematics for Visual Computing, Part 2, Ch. 12)}
\end{itemize}

\subsection{Source coverage}\label{appendix-db3aef8f2fd6-source-coverage}

\hypertarget{appendix-db3aef8f2fd6-source-coverage}{}

Supplied lesson parts: 981.

\subsection{Future keywords}\label{appendix-db3aef8f2fd6-future-keywords}

\hypertarget{appendix-db3aef8f2fd6-future-keywords}{}

endpoint tangents, C1 continuity, parameter scaling
